\begin{explanationbox}
\textbf{先看一个三边都能心算的实际三角形。}
\begin{center}\begin{tikzpicture}[font=\small,>=stealth]
\draw[black!30,rounded corners=2pt] (0,0) rectangle (14.4,9.2);
\node[anchor=west,font=\small\bfseries] at (.5,8.7) {同一个 $3$-$4$-$5$ 三角形：底高法与半角公式都算出面积 $6$};
\begin{scope}[shift={(4.5,4.2)},scale=.85]
\draw[black!35] (-3.1,0)--(3.1,0);
\fill[orange!12] (-2,0)--(2,3)--(2,0)--cycle;
\draw[blue,very thick,domain=-2.2:4,samples=101,variable=\t] plot ({sqrt(1+\t*\t/3)},\t);
\draw[blue,very thick,domain=-2.2:4,samples=101,variable=\t] plot ({-sqrt(1+\t*\t/3)},\t);
\draw[orange!85!black,very thick] (-2,0)--(2,3)--(2,0)--cycle;
\draw (1.7,0)--(1.7,.3)--(2,.3);
\fill (-2,0) circle (2.5pt) node[below left=6pt] {$F_1(-2,0)$};
\fill (2,0) circle (2.5pt) node[below right=6pt] {$F_2(2,0)$};
\fill (2,3) circle (2.5pt) node[above left=8pt] {$P(2,3)$};
\node[above=6pt,rotate=36.869898] at (0,1.5) {$m=5$};
\node[right=7pt] at (2,1.1) {$n=h=3$};
\draw (2,3) ++(-143.130102:.65) arc[start angle=-143.130102,end angle=-90,radius=.65];
\node at (1.1,2.7) {$\theta$};
\node at (0,-.85) {$2c=4$};
\end{scope}
\node[anchor=north west,align=left] at (9,7.7) {$a=1,\ b^2=3,\ c=2$\\[12pt]$m-n=5-3=2a$\\[12pt]$\displaystyle S=\frac12\cdot4\cdot3=6$\\[12pt]$\displaystyle\cot\frac\theta2=2$\\[12pt]$S=3\cdot2=6$};
\node at (7.2,.65) {$P$ 在 $x^2-y^2/3=1$ 上；顶角在 $P$，竖直焦半径同时就是高。};
\end{tikzpicture}
\end{center}
图中 $P=(2,3)$ 在 $x^2-y^2/3=1$ 上。焦距为 $4$，高为 $3$，
所以面积 $S=6$；两条焦半径为 $5$ 和 $3$，差为 $2=2a$。
此时 $\cos\theta=3/5$、$\sin\theta=4/5$，故
\[
 \cot\frac\theta2=\frac{1+\cos\theta}{\sin\theta}=2,
 \qquad b^2\cot\frac\theta2=3\cdot2=6.
\]
\textbf{推广一步：}焦点连线作底，始终有 $S=\tfrac12(2c)h=ch$。
横向标准式为 $c|y|$，纵向标准式为 $c|x|$。
半角式与底高式计算的是\textbf{同一个三角形}，可以互相回验。
\end{explanationbox}
\clearpage
\begin{explanationbox}
\textbf{把椭圆和双曲线放在同一个尺度下看。}
\begin{center}\begin{tikzpicture}[font=\small,>=stealth]
\draw[black!30,rounded corners=2pt] (0,0) rectangle (14.4,9.7);
\node[anchor=west,font=\small\bfseries] at (.5,9.2) {同一焦点三角形：椭圆固定两腰之和，双曲线固定两腰之差};
\node at (3.6,8.35) {椭圆：$a=4,\ b^2=12$};
\node at (10.8,8.35) {双曲线：$a=1,\ b^2=3$};
\foreach \cx in {3.6,10.8}{
\begin{scope}[shift={(\cx,4.8)},scale=.55]
\fill[orange!12] (-2,0)--(2,3)--(2,0)--cycle;
\draw[orange!85!black,very thick] (-2,0)--(2,3)--(2,0)--cycle;
\fill (-2,0) circle (2.5pt) node[below=7pt] {$F_1$};
\fill (2,0) circle (2.5pt) node[below=7pt] {$F_2$};
\fill (2,3) circle (2.5pt);
\node[above left=5pt] at (0,1.5) {$5$};
\node[right=6pt] at (2,1.2) {$3$};
\draw (2,3) ++(-143.130102:.65) arc[start angle=-143.130102,end angle=-90,radius=.65];
\end{scope}}
\begin{scope}[shift={(3.6,4.8)},scale=.55]
\draw[blue,very thick] (0,0) ellipse (4 and {sqrt(12)});
\node at (2.7,3.8) {$P$};
\end{scope}
\begin{scope}[shift={(10.8,4.8)},scale=.55]
\draw[blue,very thick,domain=-3.5:4.3,samples=101,variable=\t] plot ({sqrt(1+\t*\t/3)},\t);
\draw[blue,very thick,domain=-3.5:4.3,samples=101,variable=\t] plot ({-sqrt(1+\t*\t/3)},\t);
\node at (1.2,3.8) {$P$};
\end{scope}
\node at (3.6,2.25) {$5+3=8=2a$};
\node at (10.8,2.25) {$5-3=2=2a$};
\node at (3.6,1.5) {$S=12\tan(\theta/2)=12\cdot\tfrac12=6$};
\node at (10.8,1.5) {$S=3\cot(\theta/2)=3\cdot2=6$};
\node at (7.2,.6) {同尺度、同焦点 $c=2$、同一 $P=(2,3)$、同一个 $\theta$；两条曲线的 $a,b$ 不同。};
\end{tikzpicture}\end{center}

两图的焦点、顶点 $P$ 和三条边完全相同，故顶角与面积相同；
但椭圆用 $5+3=2a$，双曲线用 $5-3=2a$，相应的 $b^2$ 也不同。
\textbf{不能固定同一个 $b^2$，再机械互换 $\tan$ 与 $\cot$。}

从数字到一般式：椭圆固定 $m+n$，对应
$b^2=mn\cos^2(\theta/2)$；双曲线固定 $|m-n|$，对应
$b^2=mn\sin^2(\theta/2)$。再用
$S=mn\sin(\theta/2)\cos(\theta/2)$，就分别得到正切与余切。

\textbf{看变化方向：}例5后的三态图会显示，双曲线上高度越大，面积越大，顶角反而越小。
由 $\cot(\theta/2)=ch/b^2$，任给 $0<\theta<\pi$，取
\[
 h=\frac{b^2}{c}\cot\frac\theta2>0,\qquad
 P=\left(a\sqrt{1+h^2/b^2},h\right),
\]
即可在横向双曲线上实现该角；其他取向由坐标变换得到。
\end{explanationbox}
\clearpage

\clearpage
\begin{explanationbox}
\textbf{为什么双曲线适合用旁切圆？——只追踪“绳子差”。}
椭圆给的是两根焦半径的\textbf{和}，用内切圆，切线长相等会把这个和整理出来；
双曲线给的是两根焦半径的\textbf{差}，改用旁切圆，沿直线延长两边，
切线长相等便把“差”整理成 $PT=c-a$。

\textbf{看图的顺序：}
先找距离 $P$ 较远的焦点 $F_1$、较近的焦点 $F_2$；
再找\textbf{贴着 $PF_2$ 线段}的旁切圆，不是贴着 $F_1F_2$ 线段的那个。
圆还要同时贴住 $F_1P$ 直线越过 $P$ 的延长部分、
以及 $F_1F_2$ 直线越过 $F_2$ 的延长部分。
三个切点分别记为 $T,U,V$。三个切点和法线请参看方法五中的精确 TikZ 图。

\textbf{一句话理解面积：}
旁切圆心到三条边所在直线的垂距都等于 $r$，把面积拆成两个三角形面积之和、再减去一个三角形面积，恰好得到
\[
 S=\frac r2\bigl(PF_1+F_1F_2-PF_2\bigr)
   =\frac r2(2a+2c)=(a+c)r.
\]
而 $PT=c-a$；外角平分线给 $r/ PT=\cot(\theta/2)$。
于是面积等于 $(c+a)(c-a)\cot(\theta/2)$，
最后才用 $c^2-a^2=b^2$，\textbf{没有焦半径平方展开}。
\end{explanationbox}
